% Generic top-conference method chapter template.
% Replace bracketed placeholders with paper-specific objects.

\section{[Method Title]}\label{sec:method}

Given [input object] and [task condition], our method runs [number] stages before producing [output object]. The first stage constructs [object one], which captures [phenomenon or constraint]. The second stage turns it into [object two], a [continuous/discrete/probabilistic] representation that can be optimized locally. The final stage builds [object three], a reusable plan that drives [prediction / inference / decoding / decision]. This section defines the three objects and shows how they fit into one pipeline.

\subsection{[Stage 1: Object or Signal]}\label{sec:stage-one}

[Phenomenon] creates the first difficulty: [why direct use is insufficient]. We therefore introduce [object one] as the quantitative handle for this subsection. For each [instance] $i$, let [symbol declarations]. We define
\begin{equation}\label{eq:stage-one-score}
    s_i(x)=\operatorname{Score}\bigl([local\ feature], [reference\ set]\bigr),
\end{equation}
where [short declaration of free symbols]. A large value of $s_i(x)$ indicates [local interpretation], while a small value indicates [contrast].

\subsection{[Stage 2: Structured Representation]}\label{sec:stage-two}

The signal $s_i$ is [coarse/noisy/discrete], but the downstream procedure requires [structured property]. We convert it into [object two] by combining [source term] with [regularizer or constraint]. Let $g_i(x)$ be [structure map or auxiliary signal]. The representation is defined by
\begin{equation}\label{eq:stage-two-object}
    z_i^{\star}
    =
    \operatorname*{arg\,min}_{z}
    \left[
        \sum_x \ell(z(x), s_i(x))
        +
        \lambda \sum_{(x,y)\in E} w_i(x,y)(z(x)-z(y))^2
    \right].
\end{equation}
The first term keeps the representation faithful to the signal, and the second term prevents [undesired behavior] across [boundary or constraint].

\begin{algorithm}[tb]
\small
\caption{[Object Two] Construction}
\label{alg:stage-two}
\begin{flushleft}
\hspace*{0.02in} \textbf{Input:}
Signal $s_i$;
structure map $g_i$;
regularization weight $\lambda$;
iteration count $T$.\\
\hspace*{0.02in} \textbf{Output:}
Structured representation $z_i^{\star}$.
\end{flushleft}
\begin{algorithmic}[1]
    \State $z_i^{(0)} \gets s_i$
    \For{$t=0$ \textbf{to} $T-1$}
        \State $z_i^{(t+1)} \gets \operatorname{Update}(z_i^{(t)}, s_i, g_i, \lambda)$ as per Eq.~\ref{eq:stage-two-object}
    \EndFor
    \State $z_i^{\star} \gets z_i^{(T)}$
    \State \textbf{return} $z_i^{\star}$
\end{algorithmic}
\end{algorithm}

\paragraph{Complexity.}
Let $n$ be the number of local units and $T$ the iteration count. Each update touches [neighborhood size] neighbors per unit, so constructing $z_i^{\star}$ costs $O(Tn)$ time and $O(n)$ space.

\subsection{[Stage 3: Plan or Allocation]}\label{sec:stage-three}

The representation $z_i^{\star}$ is [continuous/global/local], but the final system needs [discrete action or allocation]. We therefore build [object three] once per input and reuse it for [scope]. Let $m_i$ be the mass of the representation,
\begin{equation}\label{eq:stage-three-budget}
    m_i=\frac{1}{n}\sum_x z_i^{\star}(x),
    \qquad
    b_i=b_{\min}+(b_{\max}-b_{\min})m_i.
\end{equation}
The budget increases only when the representation carries broad evidence, so easy instances remain cheap while difficult instances keep more information.

\begin{algorithm}[tb]
\small
\caption{End-to-End [Method Name] Pipeline}
\label{alg:e2e}
\begin{flushleft}
\hspace*{0.02in} \textbf{Input:}
Instances $\{I_i\}_{i=1}^{N}$;
model or solver $\mathcal{M}$;
stage-one parameters;
stage-two parameters;
budget bounds $(b_{\min}, b_{\max})$.\\
\hspace*{0.02in} \textbf{Output:}
Outputs $\{y_i\}_{i=1}^{N}$.
\end{flushleft}
\begin{algorithmic}[1]
    \For{$i=1$ \textbf{to} $N$}
        \State $s_i \gets$ construct stage-one signal for $I_i$
        \State $z_i^{\star} \gets$ \textsc{StageTwo}$(s_i)$
        \State $b_i \gets$ adaptive budget from Eq.~\ref{eq:stage-three-budget}
        \State $\mathcal{P}_i \gets$ build plan from $(z_i^{\star}, b_i)$
        \State $y_i \gets \mathcal{M}(I_i, \mathcal{P}_i)$
    \EndFor
    \State \textbf{return} $\{y_i\}_{i=1}^{N}$
\end{algorithmic}
\end{algorithm}

\paragraph{End-to-end complexity.}
The stage-one signal costs [cost]. Stage two costs $O(Tn)$ per instance. Stage three costs [cost] to sort, allocate, or gather. The total cost is [dominant bound], and the reusable plan amortizes its construction over [reuse scope].
